\section{Implementation and Datasets} 
\label{append:implementation-datasets}
We use a workstation running Ubuntu 24.04 LTS with an NVIDIA GeForce RTX 4090 GPU (CUDA 12.6; NVIDIA driver 560.35.03) as the primary client environment to access remote compute resources and run supporting experiments. 4 LLMs, DeepSeek-V3.2~\cite{deepseek_v32_release_2025}, Qwen-Plus~\cite{qwen_plus_modelstudio_2026}, DeepSeek-R1-8B~\cite{deepseek_r1_distill_llama8b_2025}, and Qwen3-4B~\cite{qwen3_4b_instruct_2507} are used as the base models in the evaluation. We access the first two models through their public APIs. The later two are self-hosted on Jetstream2~\cite{jetstream2025llminference} using two H100s. DeepSeek-V3.2-Release is primarily used as the base model to build \core.  
The implementations for vulnerability validation differ across vulnerabilities; we describe the corresponding setups and procedures in the relevant vulnerability sections. 


\noindent \textbf{Specification Dataset.} To ensure comparability with prior works, we construct a dataset from the same 3GPP specifications that govern control-plane signaling for NAS and RRC, which are critical for the security and mobility of operational 5G and 4G networks. In addition, to showcase that \lm can be applied to other 3GPP standards, specifications for 5G system architecture and SIP/SDP protocols are also included.
Thus, the dataset includes: \textbf{\textit{TS 38.331}} version 17.0.0 Release 17~\cite{3GPP_TS_38331}, \textbf{\textit{TS 24.501}} version 16.8.0 Release 16 \cite{3GPP_TS_24501}, \textbf{\textit{TS 24.301}} version 16.8.0 Release 16 \cite{3GPP_TS_24301}, \textbf{\textit{TS 23.501}} version 20.0.0 Release 20 \cite{3GPP_TS_23501}, and \textbf{\textit{TS 24.229}} version 19.5.0 Release 19 \cite{ETSI-TS-124229}.  